\section{Inverse systems}
\label{section-inverse-systems}

\noindent
Let $\mathcal{C}$ be a category.
In Categories, Section \ref{categories-section-posets-limits}
we defined the notion of an inverse system over a preordered set
(with values in the category $\mathcal{C}$).
If the preordered set is $\mathbf{N} = \{1, 2, 3, \ldots\}$
with the usual ordering such an inverse system over $\mathbf{N}$
is often simply called an {\it inverse system}. It consists quite simply
of a pair $(M_i, f_{ii'})$ where each $M_i$, $i \in \mathbf{N}$
is an object of $\mathcal{C}$, and for each $i > i'$, $i, i' \in \mathbf{N}$
a morphism $f_{ii'} : M_i \to M_{i'}$ such that moreover
$f_{i'i''} \circ f_{ii'} = f_{ii''}$ whenever this makes sense.
It is clear that in fact it suffices to give the morphisms
$M_2 \to M_1$, $M_3 \to M_2$, and so on. Hence an inverse system
is frequently pictured as follows
$$
M_1 \xleftarrow{\varphi_2} M_2 \xleftarrow{\varphi_3} M_3 \leftarrow \ldots
$$
Moreover, we often omit the transition maps $\varphi_i$ from the notation
and we simply say ``let $(M_i)$ be an inverse system''.

\medskip\noindent
The collection of all inverse systems with values in
$\mathcal{C}$ forms a category with the obvious notion of morphism.

\begin{lemma}
\label{lemma-inverse-systems-abelian}
Let $\mathcal{C}$ be a category.
\begin{enumerate}
\item If $\mathcal{C}$ is an additive category, then the category
of inverse systems with values in $\mathcal{C}$ is an additive category.
\item If $\mathcal{C}$ is an abelian category, then the category
of inverse systems with values in $\mathcal{C}$ is an abelian category.
A sequence $(K_i) \to (L_i) \to (M_i)$ of inverse systems
is exact if and only if each $K_i \to L_i \to N_i$ is exact.
\end{enumerate}
\end{lemma}

\begin{proof}
Omitted.
\end{proof}

\noindent
The limit (see Categories, Section \ref{categories-section-posets-limits})
of such an inverse system is denoted $\lim M_i$, or $\lim_i M_i$.
If $\mathcal{C}$ is the category of abelian groups (or sets), then the limit
always exists and in fact can be described as follows
$$
\lim_i M_i
=
\{(x_i) \in \prod M_i \mid \varphi_i(x_i) = x_{i - 1}, \ i = 2, 3, \ldots\}
$$
see Categories, Section \ref{categories-section-limit-sets}.
However, given a short exact sequence
$$
0 \to (A_i) \to (B_i) \to (C_i) \to 0
$$
of inverse systems of abelian groups it is not always the case
that the associated system of limits is exact. In order to discuss
this further we introduce the following notion.

\begin{definition}
\label{definition-Mittag-Leffler}
Let $\mathcal{C}$ be an abelian category.
We say the inverse system $(A_i)$
satisfies the {\it Mittag-Leffler condition}, or for short
is {\it ML}, if for every $i$ there exists a $c = c(i) \geq i$
such that
$$
\Im(A_k \to A_i) = \Im(A_c \to A_i)
$$
for all $k \geq c$.
\end{definition}

\noindent
It turns out that the Mittag-Leffler condition is good enough to ensure
that the $\lim$-functor is exact, provided one works within
the abelian category of abelian groups, modules over a ring, etc.
It is shown in a paper by A.\ Neeman (see \cite{Neeman-Counterexample})
that this condition is not strong enough in an abelian
category having AB4* (having exact products).

\begin{lemma}
\label{lemma-Mittag-Leffler}
Let
$$
0 \to (A_i) \to (B_i) \to (C_i) \to 0
$$
be a short exact sequence of inverse systems of abelian groups.
\begin{enumerate}
\item In any case the sequence
$$
0 \to \lim_i A_i \to \lim_i B_i \to \lim_i C_i
$$
is exact.
\item If $(B_i)$ is ML, then also $(C_i)$ is ML.
\item If $(A_i)$ is ML, then
$$
0 \to \lim_i A_i \to \lim_i B_i \to \lim_i C_i \to 0
$$
is exact.
\end{enumerate}
\end{lemma}

\begin{proof}
Nice exercise. See
Algebra, Lemma \ref{algebra-lemma-Mittag-Leffler} for part (3).
\end{proof}

\begin{lemma}
\label{lemma-apply-Mittag-Leffler}
Let
$$
(A_i) \to (B_i) \to (C_i) \to (D_i)
$$
be an exact sequence of inverse systems of abelian groups. If the
system $(A_i)$ is ML, then the sequence
$$
\lim_i B_i \to \lim_i C_i \to \lim_i D_i
$$
is exact.
\end{lemma}

\begin{proof}
Let $Z_i = \Ker(C_i \to D_i)$ and $I_i = \Im(A_i \to B_i)$.
Then $\lim Z_i = \Ker(\lim C_i \to \lim D_i)$ and
we get a short exact sequence of systems
$$
0 \to (I_i) \to (B_i) \to (Z_i) \to 0
$$
Moreover, by
Lemma \ref{lemma-Mittag-Leffler}
we see that $(I_i)$ has (ML), thus another application of
Lemma \ref{lemma-Mittag-Leffler}
shows that $\lim B_i \to \lim Z_i$ is surjective which
proves the lemma.
\end{proof}

\noindent
The following characterization of essentially constant inverse systems
shows in particular that they have ML.

\begin{lemma}
\label{lemma-essentially-constant}
Let $\mathcal{A}$ be an abelian category.
Let $(A_i)$ be an inverse system in $\mathcal{A}$ with limit $A = \lim A_i$.
Then $(A_i)$ is essentially constant (see
Categories, Definition
\ref{categories-definition-essentially-constant-diagram})
if and only if there exists an $i$ and for all $j \geq i$ a direct sum
decomposition $A_j = A \oplus Z_j$ such that
(a) the maps $A_{j'} \to A_j$ are compatible with the direct sum
decompositions, (b) for all $j$ there exists some $j' \geq j$ such that
$Z_{j'} \to Z_j$ is zero.
\end{lemma}

\begin{proof}
Assume $(A_i)$ is essentially constant. Then there exists an $i$ and
a morphism $A_i \to A$ such that $A \to A_i \to A$ is the identity and
for all $j \geq i$ there exists
a $j' \geq j$ such that $A_{j'} \to A_j$ factors as
$A_{j'} \to A_i \to A \to A_j$ (the last map comes from $A = \lim A_i$).
Hence setting $Z_j = \Ker(A_j \to A)$ for all $j \geq i$ works.
Proof of the converse is omitted.
\end{proof}

\noindent
We will improve on the following lemma in
More on Algebra, Lemma \ref{more-algebra-lemma-Mittag-Leffler}.

\begin{lemma}
\label{lemma-exact-sequence-ML}
Let
$$
0 \to (A_i) \to (B_i) \to (C_i) \to 0
$$
be an exact sequence of inverse systems of abelian groups.
If $(C_i)$ is essentially constant, then $(A_i)$ has ML
if and only if $(B_i)$ has ML.
\end{lemma}

\begin{proof}
After renumbering we may assume that $C_i = C \oplus Z_i$ compatible with
transition maps and that for all $i$ there exists an $i' \geq i$ such that
$Z_{i'} \to Z_i$ is zero, see Lemma \ref{lemma-essentially-constant}.

\medskip\noindent
First, assume $C = 0$, i.e., we have $C_i = Z_i$.
In this case choose $1 = n_1 < n_2 < n_3 < \ldots$
such that $Z_{n_{i + 1}} \to Z_{n_i}$ is zero.
Then $B_{n_{i + 1}} \to B_{n_i}$ factors through $A_{n_i} \subset B_{n_i}$.
It follows that for $j \geq i + 1$ we have
$$
\Im(A_{n_j} \to A_{n_i}) \subset \Im(B_{n_j} \to B_{n_i}) \subset
\Im(A_{n_{j - 1}} \to A_{n_i})
$$
as subsets of $A_{n_i}$.
Thus the images $\Im(A_{n_j} \to A_{n_i})$ stabilize for $j \geq i + 1$
if and only if the same is true for the images $\Im(B_{n_j} \to B_{n_i})$.
The equivalence follows from this (small detail omitted).

\medskip\noindent
If $C \not = 0$, denote $B'_i \subset B_i$ the inverse image
of $C$ by the map $B_i \to C \oplus Z_i$. Then by the previous
paragraph we see that $(B'_i)$ has ML if and only if $(B_i)$
has ML. Thus we may replace $(B_i)$ by $(B'_i)$.
In this case we have exact sequences $0 \to A_i \to B_i \to C \to 0$
for all $i$. It follows that
$0 \to \Im(A_j \to A_i) \to \Im(B_j \to B_i) \to C \to 0$
is short exact for all $j \geq i$. Hence the images
$\Im(A_j \to A_i)$ stabilize for $j \geq i$
if and only if the same is true for $\Im(B_j \to B_i)$
as desired.
\end{proof}

\noindent
The ``correct'' version of the following lemma is
More on Algebra, Lemma \ref{more-algebra-lemma-apply-Mittag-Leffler-again}.

\begin{lemma}
\label{lemma-apply-Mittag-Leffler-again}
Let
$$
(A^{-2}_i \to A^{-1}_i \to A^0_i \to A^1_i)
$$
be an inverse system of complexes of abelian groups and denote
$A^{-2} \to A^{-1} \to A^0 \to A^1$ its limit. Denote
$(H_i^{-1})$, $(H_i^0)$ the inverse systems of cohomologies, and
denote $H^{-1}$, $H^0$ the cohomologies of $A^{-2} \to A^{-1} \to A^0 \to A^1$.
If $(A^{-2}_i)$ and $(A^{-1}_i)$ are ML and
$(H^{-1}_i)$ is essentially constant, then
$H^0 = \lim H_i^0$.
\end{lemma}

\begin{proof}
Let $Z^j_i = \Ker(A^j_i \to A^{j + 1}_i)$ and
$I^j_i = \Im(A^{j - 1}_i \to A^j_i)$.
Note that $\lim Z^0_i = \Ker(\lim A^0_i \to \lim A^1_i)$ as
taking kernels commutes with limits.
The systems $(I^{-1}_i)$ and $(I^0_i)$ have ML as quotients of
the systems $(A^{-2}_i)$ and $(A^{-1}_i)$, see
Lemma \ref{lemma-Mittag-Leffler}.
Thus an exact sequence
$$
0 \to (I^{-1}_i) \to (Z^{-1}_i) \to (H^{-1}_i) \to 0
$$
of inverse systems where $(I^{-1}_i)$ has ML
and where $(H^{-1}_i)$ is essentially constant by assumption.
Hence $(Z^{-1}_i)$ has ML by
Lemma \ref{lemma-exact-sequence-ML}.
The exact sequence
$$
0 \to (Z^{-1}_i) \to (A^{-1}_i) \to (I^0_i) \to 0
$$
and an application of
Lemma \ref{lemma-Mittag-Leffler}
shows that $\lim A^{-1}_i \to \lim I^0_i$ is surjective.
Finally, the exact sequence
$$
0 \to (I^0_i) \to (Z^0_i) \to (H^0_i) \to 0
$$
and
Lemma \ref{lemma-Mittag-Leffler}
show that $\lim I^0_i \to \lim Z^0_i \to \lim H^0_i \to 0$
is exact. Putting everything together we win.
\end{proof}

\noindent
Sometimes we need a version of the lemma above where we take limits
over big ordinals.

\begin{lemma}
\label{lemma-ML-over-ordinals}
Let $\alpha$ be an ordinal. Let $K_\beta^\bullet$, $\beta < \alpha$
be an inverse system of complexes of abelian groups over $\alpha$. If
for all $\beta < \alpha$ the complex $K_\beta^\bullet$ is acyclic and
the map
$$
K^n_\beta \longrightarrow \lim_{\gamma < \beta} K^n_\gamma
$$
is surjective, then the complex
$\lim_{\beta < \alpha} K_\beta^\bullet$ is acyclic.
\end{lemma}

\begin{proof}
By transfinite induction we prove this holds for every ordinal
$\alpha$ and every system as in the lemma. In particular, whilst
proving the result for $\alpha$ we may assume the complexes
$\lim_{\gamma < \beta} K^n_\gamma$ are acyclic.

\medskip\noindent
Let $x \in \lim_{\beta < \alpha} K^0_\beta$ with $\text{d}(x) = 0$.
We will find a $y \in \lim_{\beta < \alpha} K^{-1}_\beta$
with $\text{d}(y) = x$.
Write $x = (x_\beta)$ where $x_\beta \in K_\beta^0$ is the
image of $x$ for $\beta < \alpha$. We will construct $y = (y_\beta)$
by transfinite recursion.

\medskip\noindent
For $\beta = 0$ let $y_0 \in K_0^{-1}$
be any element with $\text{d}(y_0) = x_0$.

\medskip\noindent
For $\beta = \gamma + 1$ a successor, we have to find an element $y_\beta$
which maps both to $y_\gamma$ by the transition map
$f : K^\bullet_\beta \to K^\bullet_\gamma$ and to $x_\beta$ under the
differential. As a first approximation we choose $y'_\beta$ with
$\text{d}(y'_\beta) = x_\beta$. Then the difference $y_\gamma - f(y'_\beta)$
is in the kernel of the differential, hence equal to $\text{d}(z_\gamma)$
for some $z_\gamma \in K^{-2}_\gamma$.
By assumption, the map $f^{-2} : K^{-2}_\beta \to K^{-2}_\gamma$
is surjective. Hence we write $z_\gamma = f(z_\beta)$
and change $y'_\beta$ into $y_\beta = y'_\beta + \text{d}(z_\beta)$
which works.

\medskip\noindent
If $\beta$ is a limit ordinal, then we have the element
$(y_\gamma)_{\gamma < \beta}$ in $\lim_{\gamma < \beta} K^{-1}_\gamma$
whose differential is the image of $x_\beta$. Thus we can argue in exactly
the same manner as above using the termwise surjective map of complexes
$f : K_\beta^\bullet \to \lim_{\gamma < \beta} K_\gamma^\bullet$
and the fact (see first paragraph of proof) that we may assume
$\lim_{\gamma < \beta} K_\gamma^\bullet$ is acyclic by induction.
\end{proof}


